\documentclass[11pt,a4paper]{article}

\usepackage[margin=1in]{geometry}
\usepackage{xcolor}
\usepackage{listings}
\usepackage{amsmath}
\usepackage{booktabs}
\usepackage{enumitem}
\usepackage{pgfplots}
\usepackage[hidelinks]{hyperref}
\pgfplotsset{compat=1.18}

\definecolor{codebg}{gray}{0.96}
\definecolor{plaincol}{RGB}{185,168,148}
\definecolor{tunedcol}{RGB}{122,75,42}
\definecolor{decided}{RGB}{0,110,60}
\definecolor{open}{RGB}{170,60,0}

\lstdefinestyle{ml}{%
  language=[Objective]Caml,
  basicstyle=\ttfamily\small,
  backgroundcolor=\color{codebg},
  keywordstyle=\color{blue!60!black},
  commentstyle=\itshape\color{black!45},
  columns=fullflexible,
  breaklines=true,
  showstringspaces=false,
  frame=none,
  xleftmargin=1em,
  framexleftmargin=1em,
  aboveskip=0.6em, belowskip=0.6em
}

\pgfplotsset{
  tatamibar/.style={
    ybar, width=\textwidth, height=6.2cm,
    bar width=11pt, ymin=0,
    xtick=data, ylabel={ms},
    enlarge x limits=0.22,
    tick label style={font=\small},
    label style={font=\small},
    legend style={at={(0.5,1.03)}, anchor=south, legend columns=-1, font=\small, draw=none},
    nodes near coords={\pgfmathprintnumber[fixed,precision=0]{\pgfplotspointmeta}},
    every node near coord/.append style={font=\tiny, color=black!55, inner sep=1pt},
    ymajorgrids, grid style={black!12},
  }
}

\newcommand{\squeeze}[1]{\textcolor{open}{\textbf{#1}}}

\title{Phase 3 --- what the signature costs and what it buys}
\author{Measurements, not proposals}
\date{}

\begin{document}
\maketitle

\begin{abstract}
\noindent
Two kernels answer the same query against the same database. \texttt{plain} has never
read the generated signature and holds every value as text. \texttt{tuned} has read
\texttt{orders.mli} and holds each column as a dense array of the type the signature
names, with a validity mask only where the word \texttt{option} appears. Both send
identical SQL, receive identical rows, and stream them rather than materialising an
intermediate list. The only variable is what the rows become in memory.

The headline is \textbf{2.2$\times$}, not the order of magnitude a first estimate
suggested, and the sections below say where the rest went. \textbf{That is accepted as
sufficient}, for the reasons in \S\ref{sec:enough}. Every figure is emitted by
\texttt{dune exec bin/bench.exe} into \texttt{docs/phase\_3/data/} and read from there
by this document, so it can be rebuilt rather than believed.
\end{abstract}

\section{Method, and what is being held constant}

The corpus is 200{,}012 rows of \texttt{orders}, one row in three having a NULL
\texttt{note}. Timings are the median of seven runs; the first is discarded because it
pays for the connection and a cold page cache. Variance between whole runs is a few
percent, and \emph{no claim here rests on a difference smaller than that} --- an
earlier draft reported a 14\% effect that reversed sign on the next run.

Three things are deliberately held constant, because each is a way of accidentally
measuring something other than the signature.

\begin{enumerate}[leftmargin=1.4em]
\item \textbf{The SQL.} Both kernels push the predicate to Postgres. The database does
  identical work and the same bytes cross the wire. The only textual difference is that
  \texttt{select *} must stay a star for \texttt{plain}, which has no signature telling
  it what the table contains, while \texttt{tuned} expands it.
\item \textbf{Streaming.} An earlier version had \texttt{tuned} stream while
  \texttt{plain} built an intermediate list, which credited the \texttt{.mli} with a
  17\% saving that belonged to \texttt{execute\_iter}. Both stream now.
\item \textbf{Validation.} \texttt{Analysis.plan} rejects a query the signature cannot
  describe --- \texttt{qty < 4.5} against an \texttt{int} column --- before either
  kernel runs, so neither is answering a different question from the other.
\end{enumerate}

\noindent
The two kernels are checked for agreement on every measured query. A speed comparison
between kernels that return different answers would be worthless, so it is asserted
rather than assumed.

\section{Materialising a result set}

Time to turn one query's rows into each kernel's own answer.

\begin{center}
\begin{tikzpicture}
\begin{axis}[tatamibar,
  symbolic x coords={int,int+float,opt-text,all-six},
  xlabel={columns fetched}]
\addplot[fill=plaincol, draw=plaincol!70!black] table[x=layout,y=plain] {data/materialise.dat};
\addplot[fill=tunedcol, draw=tunedcol!70!black] table[x=layout,y=tuned] {data/materialise.dat};
\legend{plain, tuned}
\end{axis}
\end{tikzpicture}
\end{center}

The win tracks the numeric fraction of the columns, and that is the whole explanation.
\texttt{qty} alone is around 20\%: text arrives, and an \texttt{int array} is what the
signature says to put it in, so the parse happens once instead of being deferred into
every later question. A nullable text column is around 10\%, and \texttt{select *}
--- four of whose six columns are text or text-like --- less still.

This is the first honest disappointment. A string stays a string in both kernels.
There is no conversion for the signature to hoist, so on text the \texttt{.mli} buys
almost nothing at this stage. \emph{A corpus that is mostly strings is a corpus where
this claim is weakest}, which is precisely the risk already recorded in
\texttt{scratch.org} under ``what would actually kill this''.

\section{What each column costs}

Columns added one at a time, so the marginal cost of each is visible.

\begin{center}
\begin{tikzpicture}
\begin{axis}[tatamibar,
  symbolic x coords={qty,id,customer-id,price,sku,note},
  xlabel={column added}, ylabel={marginal ms},
  bar width=18pt]
\addplot[fill=tunedcol, draw=tunedcol!70!black] table[x=added,y=marginal] {data/columns.dat};
\end{axis}
\end{tikzpicture}
\end{center}

The first bar carries a fixed per-query cost --- the round trip and the row iteration
--- which is why it is not comparable with the rest. Of the others, the expensive one
is \texttt{note}, and it is expensive for three reasons that are all legible in its
type. \texttt{string option} is wider on the wire than an integer, it forces a validity
array to be allocated, and it puts a branch in the store path for every row.

\begin{center}\itshape
This table is predictable from the signature alone, before anything runs.
\end{center}

\noindent
That is the thesis in its literal form: \emph{the access shape is known for free,
because we generated the layout}. The signature does not make the fetch faster; it says
what the fetch will cost, and for a placement decision that is the half that matters.
It is also the one result here that does not depend on a benchmark at all.

\section{Load once, query many}

Materialisation alone cannot justify itself: if the answer is wanted once and then
discarded, building a dense array and reading it back out is strictly more work than
leaving the text alone. The arrays only repay if something runs over them more than
once --- which is a claim about the workload, not about the compiler.

Each query below runs \emph{in memory}, against a corpus already fetched.

\begin{center}
\begin{tikzpicture}
\begin{axis}[tatamibar,
  symbolic x coords={int-half-survive,int-few-equality,float-half-survive,text-compare,int-few-lessthan},
  x tick label style={rotate=18, anchor=east, font=\scriptsize},
  xlabel={}, height=6.6cm]
\addplot[fill=plaincol, draw=plaincol!70!black] table[x=label,y=plain] {data/workload.dat};
\addplot[fill=tunedcol, draw=tunedcol!70!black] table[x=label,y=tuned] {data/workload.dat};
\legend{plain, tuned}
\end{axis}
\end{tikzpicture}
\end{center}

The spread is systematic, and more informative than the average:

\begin{itemize}[leftmargin=1.4em]
\item \textbf{$\approx$13$\times$} when the predicate is numeric and few rows survive.
  \texttt{plain} calls \texttt{int\_of\_string} on every row of the corpus; \texttt{tuned}
  compares machine integers in a loop with no validity test, because \texttt{qty} is
  \texttt{int} and not \texttt{int option}.
\item \textbf{$\approx$2$\times$} when the predicate is numeric but half the table
  survives. The comparison is still 13$\times$ cheaper; what closes the gap is
  materialising the answer.
\item \textbf{$\approx$1.9$\times$} on a text predicate --- and most of that is
  \texttt{plain} walking a list to reach the column, not anything about types.
\end{itemize}

\section{Where it crosses over}
\label{sec:cross}

Total time after $k$ rounds of that five-query workload, including the one-off load.

\begin{center}
\begin{tikzpicture}
\begin{axis}[width=\textwidth, height=6.4cm,
  xlabel={rounds of the workload ($k$)}, ylabel={total ms},
  ymajorgrids, grid style={black!12}, ymin=0,
  legend style={at={(0.02,0.98)}, anchor=north west, font=\small, draw=none},
  tick label style={font=\small}, label style={font=\small}]
\addplot[plaincol, thick, mark=*, mark size=1.4] table[x=k,y=plain] {data/rounds.dat};
\addplot[tunedcol, thick, mark=*, mark size=1.4] table[x=k,y=tuned] {data/rounds.dat};
\legend{plain, tuned}
\end{axis}
\end{tikzpicture}
\end{center}

At $k=0$ the two are level: decoding costs about what \texttt{plain} saves by not
decoding. The curves separate immediately and flatten towards \textbf{2.2$\times$}. The
asymptote is not the 13$\times$ of the scan, because every query still has to build its
answer, and that allocation is paid by both.

\section{Is 2.2$\times$ enough}
\label{sec:enough}

Yes, and the reasons are worth stating rather than leaving to be inferred from the fact
that we stopped.

A factor of two is not a constant-factor rounding error. It is the difference between a
workload finishing in an afternoon and finishing overnight, and it holds asymptotically
rather than at one convenient point on the curve. Nothing about it degrades as the
corpus grows, because the plan is not a function of the corpus.

It is also the \emph{floor} of what this arrangement can show, not the ceiling, for
three separate reasons.

\begin{enumerate}[leftmargin=1.4em]
\item \textbf{The thing being computed is as thin as it can be.} One predicate, one
  projection, no aggregates, no joins. That is the least work it is possible to do over
  a materialised column and still be doing something, and it is the case least
  favourable to a typed representation. Every operation the query language gains --- a
  sum, a second predicate, a group-by --- runs over arrays that have already been paid
  for, and lands entirely on the right-hand side of the comparison.
\item \textbf{The gap that remains is identified and mechanical.} \S\ref{sec:left}
  is not a list of hopes. The survivor \texttt{int list} is measurably what collapses
  13$\times$ to 2$\times$, and a bitmask is a known replacement rather than a research
  question. We are reporting the number the current implementation earns, not the one
  it could be made to earn.
\item \textbf{The baseline has not been left slow on purpose.} \texttt{plain} streams,
  pushes its predicate down, and is charged nothing for the render that \texttt{tuned}
  pays for. Where it is still carrying avoidable cost --- \texttt{List.nth} --- that is
  written down rather than banked. A 2.2$\times$ measured against a baseline that has
  been given every fair advantage is worth more than a larger figure measured against
  one that has not.
\end{enumerate}

\noindent
And the speed was never the load-bearing claim. \S3 is: the cost of a column is
readable off its type before anything runs, and the plan that follows is a function of
the signature and not of the data. A plan that is \emph{merely} twice as fast, arrived
at with no cost model, no statistics and no profiling pass, is a stronger result than a
faster plan that needed all three --- because the former composes with a placement
decision and the latter does not. If the number were 1.0$\times$ the argument would be
in trouble. At 2.2$\times$ it is not, and chasing 13$\times$ before there is anything
to place would be optimising the part of the system that is not the question.

\section{What is left on the table}
\label{sec:left}

Everything below is measured-but-unfixed. None of it is required for the claim; all of
it is where the remaining distance between 2.2$\times$ and 13$\times$ has gone.

\begin{description}[leftmargin=0pt, style=nextline]

\item[\squeeze{The survivor list is the bottleneck.}]
\texttt{Tuned.surviving} returns an \texttt{int list} and \texttt{gather} walks it, so a
query that keeps 100{,}000 rows allocates 100{,}000 cons cells and chases them. This,
not the comparison, is what collapses 13$\times$ to 2$\times$ on unselective queries.
The replacement is a bitmask --- one bit per row in a \texttt{Bytes.t}, then a single
counting pass to size the output and a second to fill it. Two dense sequential passes
instead of a pointer chase, no per-survivor allocation, and the mask is exactly the
representation a GPU kernel would produce anyway (\S K2 of \texttt{kernels.tex}).
\texttt{plain} cannot follow: it has no dense array to mask.

\item[\squeeze{Fuse the scan with the copy.}]
With a bitmask, selective queries can skip materialising the intermediate entirely and
write survivors straight into the output array as the scan finds them. One pass rather
than three.

\item[\squeeze{\texttt{plain} is carrying avoidable cost.}]
It reaches a column with \texttt{List.nth}, which is $O(\text{column index})$ per row.
Giving it a positional fast path would shrink the reported advantage, and it should be
done before any of the above --- a baseline that has been left slow makes the other
numbers meaningless.

\item[\squeeze{\texttt{render} is charged to \texttt{tuned} and should not be.}]
Converting arrays back to text for an HTML table is work \texttt{plain} never does,
because \texttt{plain} never left text. It is a property of the destination, not of the
representation. On \texttt{select *} it was 208\,ms --- enough on its own to invert the
result --- and it is reported as its own phase for that reason.

\item[\squeeze{\texttt{Data.cell} formats floats with \texttt{Printf.sprintf "\%.2f"}.}]
Roughly four times the cost of the alternative, and it is inside the render phase above.

\item[\squeeze{Literals are pasted into the SQL text.}]
\texttt{Query.value\_to\_string} exists to display a query and does not re-escape the
quote the lexer already unescaped, so \texttt{where sku = 'a''b'} emits \texttt{= 'a'b'}
--- a syntax error, and an injection. The fix is a pgx parameter, not a better quoter.
Correctness, not performance, and it should be done first.

\item[\squeeze{The plan carries names, not layouts.}]
\texttt{Analysis.filter} and \texttt{Analysis.project} hold column names, so
\texttt{tuned} must re-derive the layout and carries unreachable \texttt{failwith}
branches for questions the plan has already answered. Holding \texttt{Schema.column}
instead would make the plan self-sufficient.

\end{description}

\section{What this does not show}

\begin{itemize}[leftmargin=1.4em]
\item \textbf{Nothing here is about a GPU.} Both kernels run on one CPU core. What has
  been established is narrower: that a plan derived from the signature alone is
  \emph{executable}, and that the representation it dictates is cheaper to query. Whether
  the same signature can decide device placement is untouched.
\item \textbf{The workload is assumed, not observed.} Load-once/query-many is the shape
  that makes shredding worth paying for, and it was chosen for that reason. A workload of
  single ad-hoc queries would show 1.0$\times$, and the honest reading of
  \S\ref{sec:cross} is that the claim depends on a property of the user, not of the
  compiler.
\item \textbf{One table, one predicate, no joins.} No \texttt{AND}/\texttt{OR}, no
  aggregates, no \texttt{ORDER BY}, no subqueries. A join would introduce the dependent
  traversal that \texttt{kernels.tex} argues is the interesting case, and it is absent
  here entirely.
\item \textbf{The corpus is synthetic and uniform.} Values are arithmetic rather than
  random, so a re-seed is reproducible --- but selectivity is therefore smooth, and
  nothing here says what happens on skew.
\end{itemize}

\end{document}
